\documentclass[a4paper,12pt]{article}

% -------------------------------------------------
% Кодировка, язык, математика
% -------------------------------------------------
\usepackage[T2A]{fontenc}
\usepackage[utf8]{inputenc}
\usepackage[russian]{babel}
\usepackage{amsmath,amssymb,mathtools}
\usepackage{helvet}
\renewcommand{\familydefault}{\sfdefault}
\usepackage{icomma}

% -------------------------------------------------
% Вёрстка
% -------------------------------------------------
\usepackage[
  a4paper,
  left=20mm,
  right=20mm,
  top=20mm,
  bottom=22mm
]{geometry}
\usepackage{microtype}
\usepackage{parskip}
\usepackage{setspace}
\onehalfspacing
\usepackage{enumitem}
\usepackage{tabularx,booktabs,longtable}
\usepackage{xcolor}
\usepackage{hyperref}
\usepackage{fancyhdr}

% -------------------------------------------------
% TikZ
% -------------------------------------------------
\usepackage{tikz}
\usetikzlibrary{calc,arrows.meta,positioning,shapes.geometric}

% -------------------------------------------------
% Цвета
% -------------------------------------------------
\definecolor{accent}{HTML}{008080}
\definecolor{darkaccent}{HTML}{004D4D}
\definecolor{textgray}{HTML}{333333}
\definecolor{softgray}{HTML}{6F7777}
\definecolor{lightaccent}{HTML}{E7F1F1}

\color{textgray}

% -------------------------------------------------
% Ссылки
% -------------------------------------------------
\hypersetup{
  colorlinks=true,
  linkcolor=darkaccent,
  urlcolor=darkaccent,
  pdftitle={Чек-лист: степени, корни, уравнения и системы},
  pdfauthor={Лёвин Артём Александрович}
}

% -------------------------------------------------
% Колонтитулы
% -------------------------------------------------
\pagestyle{fancy}
\fancyhf{}
\fancyhead[L]{\small Ярослав Верещагин}
\fancyhead[R]{\small ОГЭ по математике}
\fancyfoot[C]{\small\thepage}
\renewcommand{\headrulewidth}{0.3pt}
\renewcommand{\footrulewidth}{0pt}

% -------------------------------------------------
% Списки
% -------------------------------------------------
\setlist[itemize]{
  leftmargin=6mm,
  itemsep=2pt,
  topsep=3pt
}

\setlist[enumerate]{
  leftmargin=7mm,
  itemsep=5pt,
  topsep=4pt
}

% -------------------------------------------------
% Заголовки
% -------------------------------------------------
\usepackage{titlesec}

\titleformat{\section}
  {\Large\bfseries\color{darkaccent}}
  {}
  {0pt}
  {}

\titleformat{\subsection}
  {\large\bfseries\color{accent}}
  {}
  {0pt}
  {}

\titlespacing*{\section}{0pt}{8pt}{5pt}
\titlespacing*{\subsection}{0pt}{7pt}{3pt}

% -------------------------------------------------
% TikZ-стили блок-схем
% -------------------------------------------------
\tikzset{
  flowstart/.style={
    ellipse,
    draw=accent,
    line width=0.9pt,
    minimum width=38mm,
    minimum height=10mm,
    align=center,
    inner sep=4pt
  },
  flowprocess/.style={
    rectangle,
    rounded corners=3pt,
    draw=accent,
    line width=0.9pt,
    text width=104mm,
    minimum height=12mm,
    align=center,
    inner sep=5pt
  },
  flowdecision/.style={
    diamond,
    aspect=2.6,
    draw=accent,
    line width=0.9pt,
    text width=45mm,
    align=center,
    inner sep=2pt
  },
  flowarrow/.style={
    -{Stealth[length=3mm,width=2mm]},
    line width=0.8pt,
    draw=darkaccent
  }
}

% -------------------------------------------------
% Пользовательские команды
% -------------------------------------------------
\newcommand{\important}[1]{\textbf{\color{darkaccent}#1}}
\newcommand{\ruleline}{%
  \vspace{2mm}
  {\color{accent}\hrule height 0.5pt}
  \vspace{2mm}
}

\begin{document}

% =================================================
% ТИТУЛЬНЫЙ ЛИСТ
% =================================================
\begin{titlepage}
\thispagestyle{empty}

\begin{center}

\vspace*{24mm}

{\Large\color{softgray} Математика}

\vspace{8mm}

{\fontsize{27}{32}\selectfont\bfseries\color{darkaccent}
Чек-лист}

\vspace{4mm}

{\Large
Степени, квадратные корни,\\
уравнения и системы}

\vspace{15mm}

{\large
Подготовка к ОГЭ}

\vspace{12mm}

{\large\bfseries
Ярослав Верещагин}

\vfill

{\large
Лёвин Артём Александрович\\[2mm]
эксклюзивно для Ярослава Верещагина}

\vspace{8mm}

{\large \today}

\vspace{24mm}

\begin{tikzpicture}[remember picture,overlay]
  \draw[
    accent,
    line width=1.1pt
  ]
  ([xshift=23mm,yshift=25mm]current page.south west)
    .. controls
    ([xshift=65mm,yshift=38mm]current page.south west)
    and
    ([xshift=-65mm,yshift=18mm]current page.south east)
    ..
  ([xshift=-23mm,yshift=30mm]current page.south east);
\end{tikzpicture}

\end{center}
\end{titlepage}

% =================================================
% ЧЕК-ЛИСТ
% =================================================
\section{Основной чек-лист}

\subsection{1. Степени: что необходимо помнить}

Для одинакового основания:

\[
a^m\cdot a^n=a^{m+n},
\qquad
\frac{a^m}{a^n}=a^{m-n},
\quad a\ne0.
\]

При возведении степени в степень показатели перемножаются:

\[
(a^m)^n=a^{mn}.
\]

Например,

\[
\frac{(x^4)^3\cdot x^2}{x^9}
=
x^{12+2-9}
=
x^5,
\qquad x\ne0.
\]

\important{Проверьте ограничение.}
Если переменная стоит в знаменателе исходного выражения, значение, обращающее знаменатель в нуль, запрещено.

\ruleline

\subsection{2. Сокращение: что остаётся после деления}

Если одинаковые множители находятся в числителе и знаменателе, после сокращения остаётся единица:

\[
\frac{b^6}{b^6}=1,
\qquad b\ne0.
\]

Поэтому

\[
\frac{a^{14}b^6}{a^9b^6}
=
a^5.
\]

\important{Типичная ошибка:}
считать, что одинаковые множители после сокращения дают нуль.

Нуль возникает при вычитании:

\[
b^6-b^6=0.
\]

При делении:

\[
\frac{b^6}{b^6}=1.
\]

\subsection{3. Когда сокращать нельзя}

Сокращаются \important{множители}, а не отдельные слагаемые.

Например, преобразование

\[
\frac{15-7}{5}
\]

не позволяет ``сократить'' \(15\) и \(5\), поскольку числитель представляет собой разность.

Сначала выполняем действие:

\[
\frac{15-7}{5}
=
\frac85.
\]

Если выражение разложено на множители, сокращение возможно:

\[
\frac{5(x-2)}{5(x+1)}
=
\frac{x-2}{x+1},
\qquad x\ne-1.
\]

% =================================================
% КОРНИ И МОДУЛЬ
% =================================================
\section{Квадратный корень и модуль}

Главное правило:

\[
\boxed{\sqrt{a^2}=|a|}
\]

Поэтому

\[
\sqrt{x^{10}}
=
\sqrt{(x^5)^2}
=
|x^5|.
\]

Аналогично,

\[
\sqrt{81x^{10}y^4}
=
9|x^5|\,y^2.
\]

Обратите внимание:

\[
\sqrt{y^4}
=
\sqrt{(y^2)^2}
=
|y^2|
=
y^2,
\]

поскольку \(y^2\ge0\) при любом действительном \(y\).

\subsection{Почему модуль иногда обязателен}

Равенство

\[
\sqrt{x^2}=x
\]

верно только при \(x\ge0\).

Например, при \(x=-4\):

\[
\sqrt{x^2}
=
\sqrt{16}
=
4,
\]

поэтому корректная запись:

\[
\sqrt{x^2}=|x|.
\]

\subsection{Сворачивание полного квадрата}

Если под корнем находится квадрат выражения,

\[
\sqrt{(6x-5)^2}
=
|6x-5|.
\]

После подстановки конкретного \(x\) модуль раскрывается по знаку выражения.

Например, при \(x=\frac12\):

\[
|6x-5|
=
|3-5|
=
|-2|
=
2.
\]

\important{Ключевая проверка:}
после извлечения квадратного корня из квадрата переменного выражения сначала записывайте модуль.



% =================================================
% ЛИНЕЙНЫЕ УРАВНЕНИЯ
% =================================================
\section{Линейные уравнения}

Типовой порядок:

\begin{enumerate}
  \item раскрыть скобки;
  \item перенести слагаемые с \(x\) в одну часть;
  \item числа собрать в другой части;
  \item привести подобные;
  \item разделить обе части на коэффициент при \(x\).
\end{enumerate}

Пример:

\[
2-3(2x+2)+4x=-7.
\]

Раскрываем скобки:

\[
2-6x-6+4x=-7.
\]

Приводим подобные:

\[
-2x-4=-7.
\]

Переносим число:

\[
-2x=-3.
\]

Следовательно,

\[
x=\frac32.
\]

\important{При переносе слагаемого через знак равенства его знак меняется.}

% =================================================
% КВАДРАТНЫЕ УРАВНЕНИЯ
% =================================================
\section{Квадратные уравнения}

Полное квадратное уравнение:

\[
ax^2+bx+c=0,
\qquad a\ne0.
\]

Дискриминант:

\[
D=b^2-4ac.
\]

Если \(D>0\),

\[
x_{1,2}
=
\frac{-b\pm\sqrt D}{2a}.
\]

Если \(D=0\),

\[
x=-\frac{b}{2a}.
\]

Если \(D<0\), действительных корней нет.

\subsection{Пример}

Решим

\[
x^2-x-6=0.
\]

Здесь

\[
a=1,\qquad b=-1,\qquad c=-6.
\]

Тогда

\[
D=(-1)^2-4\cdot1\cdot(-6)
=1+24
=25.
\]

Получаем

\[
x_1=
\frac{-(-1)+5}{2}
=
3,
\]

\[
x_2=
\frac{-(-1)-5}{2}
=
-2.
\]

Ответ:

\[
\boxed{x=3;\;-2}.
\]

\important{Особое внимание знаку коэффициента \(b\).}
Если перед \(x\) стоит минус, то \(b<0\).



% =================================================
% НЕПОЛНЫЕ КВАДРАТНЫЕ УРАВНЕНИЯ
% =================================================
\section{Неполные квадратные уравнения}

\subsection{Случай 1: \(ax^2+bx=0\)}

Выносим \(x\):

\[
x(ax+b)=0.
\]

Произведение равно нулю, если хотя бы один множитель равен нулю.

Например,

\[
x^2-6x=0,
\]

\[
x(x-6)=0,
\]

откуда

\[
x=0
\qquad\text{или}\qquad
x=6.
\]

\subsection{Случай 2: \(ax^2+c=0\)}

Например,

\[
x^2-6=0.
\]

Переносим число:

\[
x^2=6.
\]

Поэтому

\[
x=\pm\sqrt6.
\]

Если получилось

\[
x^2=-6,
\]

действительных решений нет, поскольку

\[
x^2\ge0
\]

для любого действительного \(x\).

% =================================================
% СИСТЕМЫ
% =================================================
\section{Системы уравнений: метод подстановки}

Рассмотрим систему

\[
\begin{cases}
2x+y=10,\\
x-y=5.
\end{cases}
\]

Из второго уравнения выразим \(x\):

\[
x=y+5.
\]

Подставляем выражение в первое уравнение:

\[
2(y+5)+y=10.
\]

\important{Подставляемое выражение рекомендуется брать в скобки.}

Раскрываем скобки:

\[
2y+10+y=10.
\]

Отсюда

\[
3y=0,
\qquad
y=0.
\]

Теперь

\[
x=0+5=5.
\]

Ответ:

\[
\boxed{(5;\,0)}.
\]

\subsection{Если перед переменной стоит минус}

Из равенства

\[
2x-y=5
\]

сначала корректно выражаем именно \(y\):

\[
-y=5-2x.
\]

Умножаем обе части на \(-1\):

\[
y=2x-5.
\]

\important{Записи \(y\) и \(-y\) обозначают разные выражения.}
Перед подстановкой убедитесь, что выразили именно ту переменную, значение которой требуется подставить.

\subsection{Дробные ответы}

Если после решения получается, например,

\[
x=\frac{27}{7},
\]

лучше сохранить обычную дробь.

Переход к приближённой десятичной записи вида

\[
3{,}857142\ldots
\]

обычно только усложняет дальнейшие вычисления.

% =================================================
% БЛОК-СХЕМА СИСТЕМ
% =================================================
\clearpage
\thispagestyle{fancy}

\begin{center}
{\LARGE\bfseries\color{darkaccent}
Алгоритм: метод подстановки}
\end{center}

\vspace{7mm}

\begin{center}
\begin{tikzpicture}[node distance=10mm]

\node[flowstart] (start)
{Дана система двух уравнений};

\node[flowprocess, below=of start] (choose)
{Выберите уравнение, из которого удобно выразить одну переменную};

\node[flowprocess, below=of choose] (express)
{Выразите переменную полностью, внимательно обработав знаки};

\node[flowprocess, below=of express] (substitute)
{Подставьте полученное выражение во второе уравнение\\
и обязательно заключите его в скобки};

\node[flowprocess, below=of substitute] (solve)
{Решите получившееся уравнение с одной переменной};

\node[flowprocess, below=of solve] (back)
{Подставьте найденное значение обратно\\
и найдите вторую переменную};

\node[flowprocess, below=of back] (check)
{Проверьте пару подстановкой в исходные уравнения};

\node[flowstart, below=of check] (finish)
{Запишите пару значений};

\draw[flowarrow] (start) -- (choose);
\draw[flowarrow] (choose) -- (express);
\draw[flowarrow] (express) -- (substitute);
\draw[flowarrow] (substitute) -- (solve);
\draw[flowarrow] (solve) -- (back);
\draw[flowarrow] (back) -- (check);
\draw[flowarrow] (check) -- (finish);

\end{tikzpicture}
\end{center}

\clearpage

% =================================================
% ТИПИЧНЫЕ ОШИБКИ
% =================================================
\section{Что особенно важно контролировать}

\begin{enumerate}
  \item
  При извлечении корня:
  \[
  \sqrt{A^2}=|A|.
  \]

  \item
  После сокращения одинаковых ненулевых множителей остаётся \(1\).

  \item
  Сокращать отдельные слагаемые через знаки \(+\) и \(-\) нельзя.

  \item
  Перед подстановкой в систему убедитесь, что выражена сама переменная, а не величина с противоположным знаком.

  \item
  Подставляемое алгебраическое выражение заключайте в скобки.

  \item
  В квадратном уравнении сначала точно определите
  \[
  a,\quad b,\quad c,
  \]
  включая их знаки.

  \item
  Если ответ получился обычной дробью, сохраняйте её в точном виде до конца вычислений.

  \item
  В выражениях со знаменателями обязательно проверяйте условие
  \[
  \text{знаменатель}\ne0.
  \]

  \item
  Перед финальным ответом полезно выполнить короткую проверку подстановкой.
\end{enumerate}

\subsection{Мини-чек перед сдачей задачи}

\begin{itemize}
  \item Я правильно переписал условие?
  \item Я не потерял знак минус?
  \item Я поставил модуль после \(\sqrt{A^2}\)?
  \item Я сокращал только множители?
  \item Я учёл ограничения знаменателей?
  \item Я взял подставляемое выражение в скобки?
  \item Я сохранил точную обычную дробь?
  \item Ответ соответствует найденным значениям?
\end{itemize}

% =================================================
% ТРЕНИРОВОЧНЫЙ БЛОК
% =================================================
\clearpage

\section{Тренировочный блок: 5 заданий}

Решите задания последовательно. Сложность постепенно возрастает.

\begin{enumerate}

\item
Упростите выражение:
\[
\frac{(x^3)^4\cdot x^5}{x^{10}},
\qquad x\ne0.
\]

\item
Упростите выражение:
\[
\sqrt{49a^6b^4}.
\]

Запишите результат корректно для любых действительных \(a\) и \(b\).

\item
Решите уравнение:
\[
3(2x-1)-4x=9.
\]

\item
Решите квадратное уравнение:
\[
x^2-5x-14=0.
\]

\item
Решите систему методом подстановки:
\[
\begin{cases}
2x-y=7,\\
3x+2y=7.
\end{cases}
\]

\end{enumerate}

\vfill

{\small\color{softgray}
После решения проверьте прежде всего знаки, скобки при подстановке,
модуль при извлечении квадратного корня и допустимость сокращений.}

% =================================================
% ОТВЕТЫ
% =================================================
\clearpage

\section{Ответы к тренировочному блоку}

\renewcommand{\arraystretch}{1.45}

\begin{table}[ht]
\centering
\caption{Ответы}
\begin{tabularx}{\linewidth}{@{}c X@{}}
\toprule
\textbf{№} & \textbf{Ответ} \\
\midrule

1 &
\[
x^7
\]
\\

2 &
\[
7|a^3|b^2
\]
\\

3 &
\[
x=6
\]
\\

4 &
\[
x_1=7,\qquad x_2=-2
\]
\\

5 &
\[
x=3,\qquad y=-1
\]
\\

\bottomrule
\end{tabularx}
\end{table}

\vspace{8mm}

\subsection{Контрольные ориентиры}

В задании 1:

\[
(x^3)^4=x^{12},
\]

поэтому

\[
12+5-10=7.
\]

В задании 2:

\[
\sqrt{a^6}
=
\sqrt{(a^3)^2}
=
|a^3|,
\]

а

\[
\sqrt{b^4}=b^2.
\]

В задании 4:

\[
D=(-5)^2-4\cdot1\cdot(-14)
=25+56
=81.
\]

В задании 5 из первого уравнения удобно выразить

\[
y=2x-7.
\]

После подстановки:

\[
3x+2(2x-7)=7,
\]

\[
7x=21,
\qquad
x=3,
\]

\[
y=-1.
\]

\vfill

\begin{center}
{\color{darkaccent}\bfseries
Главная цель текущего этапа --- уверенно выполнять базовые преобразования
и не терять точность на знаках, модулях, дробях и подстановках.}
\end{center}

\end{document}